\section{Introduction}
\label{sec:intro}

% - We want to recover the material properties
% - Given a defined action (Via reward function) what material properties will optimize this
% - Is latent compression good
% - We use a surrogate model to have a continuous space to optimize over
% Recoding the material properties is important for making digital twins
% The topology is a problem and exposing interior is tricky
% Using fracture simulator that is non-differentiable

Recent advances in visual tracking and 3D reconstruction have made it increasingly practical to observe and quantify the physical behavior of food during manipulation, from robotic cutting~\cite{heiden2021disect, xu2023roboninja, beltran2024sliceit} to deformable object handling~\cite{shi2023robocook} and physics-aware visual synthesis~\cite{wu2025fruitninja, xie2024physgaussian}. A natural next step is to close the loop: given visual observations of food behavior, infer the underlying physical material parameters that would reproduce that behavior in simulation. Such \emph{inverse material estimation} would enable building digital twins of food items from video, calibrating simulators for robotic manipulation planning, and generating physically plausible training data for food recognition models.

This paper investigates the inverse estimation component of this pipeline. Assuming that behavioral features (such as separation distances, damage patterns, and fracture timing) have been extracted from observations, we ask: how can we efficiently recover the material parameters that produce matching behavior in a fracture simulator? We focus on this inverse step because, while visual tracking and feature extraction have mature solutions~\cite{wu2025fruitninja}, the inverse estimation problem remains open for non-differentiable fracture simulators where gradient-based inversion is impossible.

% The \emph{inverse problem} of determining which material parameters cause a simulator to reproduce a target behavior is particularly difficult for food fracture. Fracture simulators based on continuum damage mechanics~\cite{wolper2019cdmpm, wolper2020anisompm} are non-differentiable because crack nucleation and propagation introduce topological discontinuities. Each simulation takes minutes, and the response landscape is chaotic: small parameter perturbations can produce qualitatively different fracture patterns. These properties preclude gradient-based inversion and motivate surrogate-based approaches. Although recent work has demonstrated inverse material estimation for differentiable simulators using neural constitutive laws~\cite{ma2023nclaw} or Bayesian calibration from force measurements~\cite{beltran2024sliceit, heiden2021disect}, these techniques require simulator differentiability or direct force observations, neither of which is available in our setting.

We study this problem using orange peeling simulated with AnisoMPM~\cite{wolper2020anisompm}, a Material Point Method (MPM) simulator with continuum damage mechanics and anisotropic fiber reinforcement (Figure~\ref{fig:teaser}). The simulator has a $9$-dimensional material parameter space, is non-differentiable due to topological discontinuities during crack formation, takes minutes per evaluation, and exhibits chaotic sensitivity to parameters. Although recent work has demonstrated inverse estimation for differentiable simulators~\cite{ma2023nclaw, heiden2021disect} or from direct force measurements~\cite{beltran2024sliceit}, these approaches cannot be applied when the simulator itself is non-differentiable.

Our approach begins with a neural surrogate trained on $2{,}000$ forward simulations, providing a $10^6\times$ speedup that enables iterative optimization. We establish that Covariance Matrix Adaptation Evolution Strategy (CMA-ES)~\cite{hansen2001cmaes}, a standard gradient-free optimizer, can recover material parameters but requires thousands of evaluations per target, and that Proximal Policy Optimization (PPO)~\cite{schulman2017ppo}, a reinforcement learning (RL) algorithm, achieves comparable results, validating RL for inverse fracture physics.

The key limitation of both CMA-ES and per-target PPO is that they must be re-run for each new target. Since different oranges have different material properties (e.g., varying skin thickness, fiber density, and flesh firmness depending on variety and ripeness), a practical inverse system must handle arbitrary targets without retraining. We address this by training a \emph{goal-conditioned} PPO policy that conditions on the target features and learns a general inverse mapping. Once trained (${\sim}9$ minutes, one-time), this policy produces parameter estimates for any new target in $8$ evaluations (${\sim}10$\,ms). We evaluate the policy across three parameter representations: the original $9$-dimensional space and two $4$-dimensional latent spaces (VAE~\cite{kingma2014vae}, normalizing flow~\cite{dinh2017realnvp}). VAEs and normalizing flows are the two dominant generative models used for latent space optimization~\cite{gomez2018automatic, tripp2020sample, lee2025nfbo, maus2022lolbo}; we include both to evaluate whether the bijective property of flows provides advantages over the stochastic encoding of VAEs, as suggested by prior work~\cite{lee2025nfbo}. When all methods use the same surrogate evaluator, the flow latent space achieves the best recovery ($0.642$ actual), outperforming the original $9$D space by $23\%$.

% Our investigation surfaces a methodological finding of broader relevance. Initial results indicated that the raw 9D space outperformed all latent representations, apparently contradicting the hypothesis that compression aids optimization. However, an ablation revealed this to be a \emph{surrogate quality bias}: each representation space used a different neural predictor (the standalone MLP at $R^2{=}0.937$ for raw space versus embedded prediction heads at $R^2{=}0.913$ and $R^2{=}0.854$ for flow and VAE spaces, respectively). When we equalize surrogate quality by routing latent methods through the same standalone MLP (via decoding back to parameter space), the flow latent space outperforms raw space for amortized inversion by 23\% ($0.642$ vs.\ $0.523$ actual recovery). This finding highlights a systematic bias in comparisons across representation spaces: unless the same predictor is used, observed differences may reflect surrogate accuracy rather than representation quality.

Our contributions are:
\begin{itemize}
    \item A goal-conditioned RL policy that produces inverse material estimates for arbitrary targets in $8$ evaluations (${\sim}10$\,ms), achieving $0.642$ actual recovery without per-target retraining.
    \item A systematic comparison of evolutionary search and RL across the original parameter space and learned latent spaces for non-differentiable food fracture simulation.
    \item An ablation showing that evaluation design matters: latent methods using their own prediction heads perform worse than those using a shared, higher-accuracy surrogate, due to surrogate exploitation.
    \item A warm-start hybrid that initializes CMA-ES from the policy's output, achieving the best overall recovery ($0.828$, $540$ evaluations).
\end{itemize}

